%% kapa: exam=1 round=21 year=2010 period=1 form=A subject=civil_law subject_ko=민법 seq=16 number=16
%% source: Q-Net 감정평가사 기출문제 (cst00309) artlSeq=5024181 file=extracted/1008/5024181/2010년 제21회 감정평가사 1차시험 기출문제/1교시.hwp
%% conversion: kordoc 4.24.51 -> tools/gen1.py (auto); review=proofread-sonnet; flags=none
\begin{question}{16}
첨부에 관한 설명으로 옳지 않은 것은? (다툼이 있으면 판례에 의함)
\begin{choices}
  \item 甲의 동산이 乙의 동산과 섞여서 식별할 수 없는 경우, 섞인 동산 사이에 주종을 구별할 수 없다면 甲과 乙은 혼화 당시의 가액의 비율로 합성물을 공유한다.
  \item 甲의 동산을 재료로 하여 乙이 새로운 물건을 가공한 경우, 가공물의 소유권은 원칙적으로 甲에게 속한다.
  \item 甲의 동산과 乙의 동산이 부합하여 훼손하지 않으면 분리할 수 없게 된 경우, 원칙적으로 주된 동산의 소유자가 합성물의 소유권을 취득한다.
  \item 甲소유의 A 건물이 인접한 다른 건물과 합동(合棟)되어 건물로서의 독립성을 상실한 경우, A건물의 저당권자 乙은 합동된 건물의 매각대금 전부로부터 우선변제를 받을 수 있다.
  \item 甲이 권원 없이 乙소유의 토지에 식재한 수목이 토지에 부합되면, 乙이 그 소유권을 취득한다.
\end{choices}
\end{question}
\answer{4}
